\subsection{Particle-Wave Duality and Superposition}


\epigraph{
Newton thought that light was made up of particles, but then it was discovered, as we have seen here, that it behaves like a wave. Later, however (in the beginning of the twentieth century) it was found that light did indeed sometimes behave like a particle. Historically, the electron, for example, was thought to behave like a particle, and then it was found that in many respects it behaved like a wave. So it really behaves like neither. Now we have given up. We say: “It is like neither."
}
{\textit{Richard Feynman}}

Quantum means quantized, as in, coming in integers. 
You may have one atom, or a million, but you may not
have half an atom (unless you are really unlucky!).
Waves may also be quantized. Think of two children 
spinning a jumprope; the rope waves regularly, in 
one place. The children may change 
the shape of the wave(s), the period, and other
properties.

Whenever you measure a system, you disturb it. 
In classical physics this disturbance does not 
interfere with the system much. A falling ball 
or pendulum is unbothered by photons bouncing off 
and entering the eyes of physicists. However, a 
particle only slightly bigger than a photon may 
be very disturbed when we bounce a laser off of it 
to see what it's doing! When we measure a quantum 
particle, it forces the particle into an 
allowed state. The particle will remain in that 
state as long as it's not disturbed, and we can keep 
measuring if it's in that state.